\documentclass[../unravbook.tex]{subfiles}

\usepackage{parskip}
\setlength{\parskip}{10pt}
\setstretch{1.15}

\setenotez{
  list-heading = {\chapter*{#1}\markboth{#1}{#1}}
}

\begin{document}
\begin{spacing}{1.2}

\setcounter{chapter}{1}
\setcounter{endnote}{0}

\chapter{Internationalized} %% {{{
    %% end header ----------------------------
    
The UNGA took up the racial policies of the Union of South Africa during its \enquote{very first session,} after being asked by India to address the subject on behalf of its citizens.\footnote{United Nations, \enquote{Nelson Mandela International Day, 18 July,} (n.d.). \url{https://www.un.org/en/events/mandeladay/un_against_apartheid.shtml}.} From that point forward, it is no exaggeration to say that the ``South Africa problem'' of Apartheid\footnote{Ram C. Malhotra (Nepal), Rapporteur, ``Report of the Special Committee on the Policies of Apartheid of the Government of the Republic of South Africa, Annex II - Report of the delegation of the Special Committee on the Policies of Apartheid of the Government of the Republic of South Africa on the International Conference on Economic Sanctions against South Africa, London, 14-17 April 1964,'' Document S/5717, May 25, 1964, 33--65, 43. \url{https://documents.un.org/doc/undoc/gen/n64/116/04/pdf/n6411604.pdf}.(\enquote{he South African question has now become considerably more serious and direct for the independent African States as a result of the formidable programme of militarization that South Africa has undertaken in recent years. This programme introduces a new factor into the international character of the South Africa problem.})} --- the Afrikaans word for \enquote{apartness} that became the official name of a pervasive, multi-level \enquote{state doctrine} of race-based population control and segregation \parencite[138]{Malhotra-1964} --- gradually engulfed the world in nearly perpetual crisis and controversy. The UNGA quickly assumed a pivotal role in responding to the various uproars, as reflected by its thousands of declarations, resolutions, investigations, committee hearings and detailed reports, both early on and over the years. They were produced internally as well as by non-governmental organizations (NGOs) commissioned to investigate and advise on what the conditions exactly were and what precisely should be done about them \parencite{Leiss-1965}. This ferment of activity continued almost without pause until South Africa held its first national \enquote{one-person, one-vote} elections in April 1994, by which the ANC leader Nelson Mandela was elected president and the vast panoply of the preceding Afrikaner-dominated government's race-based laws and measures of societal control and enforcement was effectively revoked.\footnote{``The South African General Elections: 1994,'' ANC Parliamentary Caucus, August 27, 2019. \url{https://www.ancparliament.org.za/the-south-african-general-elections-1994/}.} 

While the UNGA's Secretary-General noted with pride on that occasion that it had remained involved \enquote{for more than four decades} and had \enquote{spearheaded the international campaign against apartheid,}\footnote{Document 209, Statement by the Spokesman for Secretary-General Boutros Boutros-Ghali applauding the election process in South Africa, UN Press Release SG/SM/5282-SAF/176, 6 May 1994, In: The United Nations and Apartheid, 1948-1004, The United Nations Blue Book Series, Vol. 1 (Department of Public Information United Nations: New York) 507} its beginnings were hardly momentous, when that initial resolution relating to the mistreatment of India's citizens\footnote{UNGA Resolution 44(1), adopted December 8, 1946, Treatment of Indians in the Union of South Africa. UNGA (1st sess. : 1946) \url{https://digitallibrary.un.org/record/209760?v=pdf}.} --- noting only that \enquote{friendly relations} had been \enquote{impaired} --- ``barely passed with the necessary two-thirds majority,'' supported only by \enquote{an extremely small scattering of nations from Europe (France, Iceland, and Norway)} along with \enquote{all the states from Africa, Asia, and the Middle East that had experienced the effects of racial discrimination from the West} and \enquote{all the Communist countries} \parencite[171]{Lauren-1988}.  But from 1946 until apartheid became fully established as a \enquote{crime against humanity} in 1998 by the Rome Statute, with what appeared to be a powerful enforcement arm (the International Criminal Court), the UNGA's active involvement against the vast internal control systems and external defense strategies that South Africa deployed to perpetuate apartheid was both far-reaching and profound. Even so, the methodologies revealed in the finalizing of India's requested initial resolution  --- framed as a straightforward bilateral dispute between two major nations of the South --- revealed important threads in the nature of the UNGA's role that continued nearly without interruption. 

The strength of the anti-apartheid movement grew steadily as tactics evolved, beginning in 1952, when India put forth a motion on apartheid specifically, sensing the opportunity to challenge racial inequalities and thereby help advance the peoples of newly-established Asian and African nations \parencite[43]{Klotz-1995}. But even more significantly, after 16 African states joined the UNGA over the course of 1960\footnote{Adriana  Ferreira, \enquote{UN Support for Decolonization Sixty Years On,} Cornell University Press, January 14, 2021. \url{https://www.cornellpress.cornell.edu/un-support-for-decolonization-sixty-years-on}. (\enquote{The so-called \enquote{Year of Africa}, 1960, marked the tipping point when it came to UN support for decolonization.})} and a brutal massacre in Sharpeville that March triggered intense outrage \parencite{Lodge-2011}, the UNGA became virtually the exclusive forum through which demands were made that a \enquote{new age} embodying \enquote{the winds of change} should be summoned to sweep across the globe.\footnote{Harold Macmillan, Prime Minister of the United Kingdom, used this phrase as the theme of his speech to Members of both Houses of the Parliament of the Union Of South Africa, Cape Town, on February 3, 1960. \url{https://web-archives.univ-pau.fr/english/TD2doc1.pdf}.} This resulted in particular from \enquote{the growth  of national consciousness in Africa} which reflected the dawn of a new, post-colonial order \parencite[52]{Jensen-2016}. The start of a new decade, epitomized by the black cover of the November 1960 issue of UNESCO's monthly magazine featuring only ``RACISM!'' in a large font, and announcing that ``all Member States'' were called upon ``to take all steps in their power to combat every form of racial discrimination, anti-Semitism, violence and hatred which may occur within their territories,''\footnote{The Executive Board of UNESCO  (Oct. 1960) The UNESCO Courier. \url{https://courier.unesco.org/en/articles/racism}.} exemplified the moment. So too did a unique opportunity emerge for the USSR-led ``eastern bloc'' of communist countries to meet that moment by throwing their geopolitical support enthusiastically in favor of this exciting opportunity for seismic change, or even upheaval. 

Across decades, this coalition held together until South African apartheid finally collapsed but, ``for most of the Afro-Asian states, the Soviet Union had been, during their freedom struggle, the enemy of the enemy, and, therefore, a strategic friend'' \parencite[140]{Bandyopadhyaya-1977}. Even so, the path was hardly smooth. South Africa's ties to Europe and North America, through its important role in gold and strategic mineral extractivism (critical to restoring financial stability and industrial productivity after World War II) and its strong anti-communist stance, caused the United States and other leading Western nations to resist succumbing to the mounting pressure in favor of South Africa's diplomatic and economic isolation.\footnote{Jeremy Green, ``The extractive foundations of Bretton Woods: gold, apartheid, and the racial politics of monetary order,'' January 8, 2026, Review of International Political Economy. \url{https://doi.org/10.1080/09692290.2025.2594477}. See also, Leiss 1965, 23. (\enquote{``The United States and the United Kingdom are the states under most pressure on this complex of issues \ldots. France is even stronger in its opposition to UN intervention.''})} But as its ``state doctrine'' of systemic racial oppression became further entrenched throughout the 1950s \parencite[43]{Klotz-1995}, its fate became the primary focus of both \enquote{East vs. West} and \enquote{North vs. South} tensions, with the East and South proving that ultimately, they held the long-term advantage. This had major implications for military training and cooperation pacts and strategic alliances, as African and Asian states saw South Africa's destabilizing tactics toward its neighbors as proof that its ongoing presence raised \enquote{grave \enquote{potentialities} for international friction that endangered the maintenance of international peace and security} \parencite[47]{Klotz-1995}. 

The resulting coalescence of priorities among newly independent states with Soviet technical and geopolitical assistance ensured that \enquote{the fight against apartheid gave form to the political project known as the Third World} \parencite[5]{Irwin-2012} --- a segment that accounted for \enquote{about one-third of UN membership.}\footnote{Walter Lacquer (ed.), \textit{A Dictionary of Politics} (The Free Press: New York), 498. The term ``Third World'' has far less usage today but a half century ago, its meaning was well-understood, as indicated by this 1971 dictionary.} While having differences of their own \parencite[26]{Leiss-1965}, their common position on South Africa allowed them to come together in a shared ``non-aligned posture'' which, in turn, led to their creation of collaborative development and trade arrangements reflecting a vision of \enquote{regionalism \ldots centred around attempts to delink from the logic of capitalist accumulation on a world scale, which was intrinsically structured to privilege the core imperialist countries and Western multinational corporations.}\footnote{Tricontinental: Institute for Social Research, \enquote{Sovereignty, Dignity, and Regionalism in the New International Order,} March 14, 2023. \url{https://thetricontinental.org/dossier-regionalism-new-international-order/}.} As \enquote{Afro-Asian countries long dominated by Europe} where ``growing nationalism was stimulated by a sense of racial grievance,'' they were now in a position to play a critical ``swing vote'' role in virtually all UNGA deliberations, while also seeking to chart a middle ground between the two superpowers which shifted the focus of their Cold War to the post-colonial agendas of their large and growing populations \parencite[]{Westad-2016}.

Thus, while the UNGA's overall involvement in the apartheid crisis continued for decades,\footnote{Enuga S. Reddy, \enquote{The Struggle against Apartheid: Lessons for Today's World,} UN Chronicle, September 1, 2007 \url{https://www.un.org/en/chronicle/article/struggle-against-apartheid-lessons-todays-world}. (``The General Assembly, the Economic and Social Council and the Commission on Human Rights have devoted thousands of meetings to the discussions on racial discrimination and adopted hundreds of resolutions. Other UN agencies, notably the International Labour Organization (ILO) and the United Nations Educational, Scientific and Cultural Organization (UNESCO), have made significant contributions to the common effort.'')} member-states living in proximity to apartheid also worked more locally to force change. In June 1960, the \enquote{Informal Permanent Machinery} became recognized as \enquote{the institutional form that African nationalism took in the post-colonial world} \parencite[52]{Irwin-2012}. Over time, various conference \enquote{groups} known by location --- some radical, others moderate --- led eventually to the Organisation of African Unity (1963--2002) and more specialized groups, such as Southern African Development Coordination Conference (SADCC) (1980--1992), challenged the apartheid regime's destabilizing tactics in the bush and against local resistance.\footnote{Leonard T. Kapungu, \enquote{Southern Africa and the Role of SADCC in Subregional Development,} in \textit{The Organization of African Unity after thirty years} (Yassin El-Ayouty, ed.) Praeger Publishers, 1994, 43--46.} Also, and especially important in the final decade of apartheid rule, South Africa was under growing pressure due to trade sanctions by major economic partners and disinvestment campaigns that \enquote{largely consisted of private pressure but also included some government action.}\footnote{Phiip I. Levy, \enquote{Sanctions on South Africa: What Did They Do?,} Yale University Economic Growth Center, February 1999, \url{http://www.econ.yale.edu/growth_pdf/cdp796.pdf}, 6--7.} Altogether, when the struggle that forced apartheid's end culminated in the early 1990s, Afrikaner elites who had stubbornly defended their racist system of rule finally found their growing inability to maintain the country's economy, given the choke points on external funding and access to foreign markets, to be unsustainable.\footnote{Ebrahim Salie, \enquote{International Pressure on Apartheid South Africa, 1980-1989}, 2017. \url{https://learning.bishops.org.za/history/wp-content/uploads/sites/130/2020/05/International-Pressure-on-Apartheid-S-A-Edited-version.pdf}} 

While the UNGA nevertheless continues to exude institutional pride in its own deep engagement,\footnote{Boutros Boutros-Ghali, ``Introduction: The United Nations and Apartheid 1948-1994,'' The United Nations Blue Books Series, Volume I, Department of Public Information --- United Nations, New York (199) (\enquote{The inauguration [on 10 May 1994 of Mr. Nelson Mandela as President of the Republic of South Africa] marked not merely the liberation of a country from racist tyranny, but the triumph of the aspiration of Africa for the total emancipation of the continent and the successful resolution of one of the gravest problems that had faced the United Nations from its very inception \ldots. Attending that ceremony on my first visit to South Africa to welcome the newborn nation on behalf of the world community, I was moved --- as an African and as the Secretary-General of an organization which had so resolutely supported the struggle against apartheid --- by the spirit of reconciliation that the new leaders had so successfully promoted.})} what has now been proven true over time is not that the UNGA saved the day, but that the long journey of challenging apartheid affected the UNGA's culture as much the organization played the decisive role in ending it. The specter of extreme racism, seen lurking for so long while many economically advanced nations were gradually forced away from tolerating it, girded anti-apartheid voices to see the hated system as a feature of their former colonial masters' capitalist mindsets.\footnote{Arun Kundnani, \textit{What Is Antiracism?: And Why It Means Anticapitalism}, Verso Books 2023. (``Apartheid could not be explained as simply the expression of age-old prejudices or hatreds: in very specific ways, it was a formation of the twentieth century. Apartheid South Africa thus posed in a particularly stark way the question of how capitalism and racism were related.'')} In effect, the global South and the UNGA underwent interrelated transformations, \enquote{both reflecting and shaping state and civil society discourse and movements} \parencite[227]{Abu-Laban-2020}, and both having profound impacts on organizational goals and shared understandings based on capitalism itself being antithetical to their post-colonial progress.\footnote{See also, Irwin 2013 (``The United Nations system was changed, I argue, not by apartheid but by the nation-state's triumph in the Black Atlantic, which forced a cross-section of elites to confront self-determination's conceptual relationship to state capacity, racial equality, and institutional interdependence.'')} In the words of a prominent leader of contemporary anti-racism movements, ``apartheid could not be explained as simply the expression of age-old prejudices or hatreds: in very specific ways, it was a formation of the twentieth century. Apartheid South Africa thus posed in a particularly stark way the question of how capitalism and racism were related.''\footnote{Arun Kundnani, \textit{What Is Antiracism?: And Why It Means Anticapitalism}, Verso Books 2023.}

Events leading up to the UNGA's December 1965 adoption of the International Convention on the Elimination of All Forms of Racial Discrimination (ICERD)\footnote{International Convention on the Elimination of All Forms of Racial Discrimination, December 21, 1965, UNGA (20th sess.). \url{https://digitallibrary.un.org/record/660429}.} offer a starting point for understanding how its efforts against apartheid have come to both define --- and distort --- its approach toward dealing with racial injustice. Article 3 of the ICERD, by which its 88 signatories (including Israel) and 182 ratifying state parties (including Palestine)\footnote{Summary of ICERD, UN Depository of Treaties, accessed, March 24, 2026, \url{https://treaties.un.org/doc/Publication/MTDSG/Volume\%20I/Chapter\%20IV/IV-2.en.pdf}.} agreed to ``particularly condemn'' apartheid, together with ``racial segregation,'\footnote{David Keane, \enquote{\enquote{Palestine v Israel} and the Collective Obligation to Condemn Apartheid under Article 3 of ICERD,} Melbourne Journal of International Law 23, no. 2 (2022): 251--75, 252. \url{https://search.informit.org/doi/10.3316/informit.027174151311897}. (``[U]nder ICERD, a collective obligation to condemn apartheid falls on all states parties as a treaty obligation under art 3.'')} indicates its focus on severe state-sanctioned race-based measures. But even before the ICERD was adopted at the end of 1965, the UNGA had already declared a global economic boycott (which was very unevenly enforced) against South Africa through Resolution 1761,\footnote{UNGA Resolution 1761 (XVII). The policies of apartheid of the Government of the Republic of South Africa (November 6, 1962), \url{https://docs.un.org/en/A/RES/1761\%28XVII\%29}.} the only time it has ever done so. More importantly, that resolution also created a \enquote{Special Committee Against Apartheid} to monitor and enforce the boycott. In 2024, the ANC-dominated government of South Africa contended that such a committee could be re-established if another nation (in this case, of course, Israel) is also found in violation of Article 3.\footnote{Report of the UN Secretary-General --- Advisory opinion of the International Court of Justice on the legal consequences arising from the policies and practices of Israel in the Occupied Palestinian Territory, including East Jerusalem, and from the illegality of the continued presence of Israel in the Occupied Palestinian Territory, December 19, 2024, 65. \url{https://docs.un.org/en/A/79/588}.} Thus, putting aside the ICERD's obligations, Article 3 may have already come with a \enquote{self-executing} trigger of a worldwide, UNGA-authorized boycott against any country charged with apartheid, even without any legal enforcement mechanism, which it lacked.

If penalties imposed on a UNGA member-state for engaging in apartheid had at least theoretically been created before the ICERD's adoption, why was it needed? Its broader focus came about not due to apartheid, but rather because at the start of the 1960s, the scourge of racism seemed resuscitated around the world --- even if it had never actually disappeared --- in multiple geographies. The onset of the American civil rights movement\footnote{Committee on the Elimination of Racial Discrimination, International Convention on the Elimination of All Forms of Racial Discrimination: 50 years of fighting racism. (2015). Office of the High Commissioner of Human Rights. \url{https://www.ohchr.org/en/treaty-bodies/cerd/international-convention-elimination-all-forms-racial-discrimination-50-years-fighting-racism}.} together with a sudden, mainly Europe-centered \enquote{swastika epidemic,} sparked interest and concern, with hostile influences arguably involved even in fanning outrage over racism in the domestic arenas of the United States.\footnote{Calder Walton, \enquote{What's Old Is New Again: Cold War Lessons for Countering Disinformation,} Texas National Security Review, Fall 2022. \url{https://tnsr.org/2022/09/whats-old-is-new-again-cold-war-lessons-for-countering-disinformation/}(\enquote{From the 1950s onwards, exploiting racial injustice in America became a stock in trade for Soviet disinformation. According to [former senior Soviet Bloc disinformation officer and defector to the United States, Ladislav] Bittman, Moscow instructed Eastern Bloc disinformation departments to make race in the United States a \enquote{top priority.} Amid race protests and the civil rights movement, they lacked no opportunities for inflaming and amplifying U.S. race relations.}} But the swastika incidents were particularly alarming because they received extensive international press coverage, focused on whether intricately coordinated antisemitic incitements and activities had not been laid rest by Nazi Germany's defeat in World War II after all.\footnote{David Keane, Addressing the Aggravated Meeting Points of Race and Religion, 6 U. Md. Law Journal of Race, Religion, Gender \& Class 367 (2006), 367-406, 371. \url{https://digitalcommons.law.umaryland.edu/rrgc/vol6/iss2/8}.} 

In fact, such incidents were what spurred requests for UNGA intervention into racism as a global problem more than anything else. The facts were indeed peculiar. From December 1959 through March 1960, swastika graffiti had suddenly appeared in noteworthy locations within western European cities, initially in West Germany, and eventually in places elsewhere, including North and South America.\footnote{``When The 1960s `Swastika Epidemic' Came to Canada,'' (n.d.), Canadian Anti-Hate Network. Retrieved March 21, 2026, from \url{https://www.antihate.ca/swastika_epidemic_canada_1960}.} Extensive media coverage and intense political debate followed. What is now known is that at least at its outset, the episode was a disinformation campaign engineered by Soviet intelligence.\footnote{Thomas Rid, \textit{Active Measures: The Secret History of Disinformation and Political Warfare}, Farar, Straus and Giroux: New York, 2020, pp. 130-31.} One recent reassessment considered it ``a strategic failure''\footnote{Weiss 2017.} because its presumed purpose --- undermining West Germany's status within NATO --- failed. But in the long run, it may have been an even greater strategic success than the Soviets could have ever imagined, by fomenting public outcry over whether western (or westernized) nations needed to address their inherently volatile racial and ethnic tensions. Thus, even as Pravda, \enquote{the official voice of Soviet communism,}\footnote{``Pravda Digital Archive,'' East View Information Services, December 13, 2022. \url{https://www.eastview.com/resources/gpa/pravda/}.} used the swastika epidemic to attack the ``racist'' intentions of West Germany's chancellor,\footnote{Remarks by Senator Dodd on Anti-Semitism, the Swastika Epidemic and Communism, Congressional Record - Senate, Vol. 106, Part 16, March 15, 1960, pp. 5205-5213, \url{https://www.congress.gov/86/crecb/1960/03/15/GPO-CRECB-1960-pt5-1-1.pdf} (\enquote{Encouraged by the reaction of the free world to the swastika incidents, the Communists propaganda apparatus had a field day. No charges were too extreme, no language too extravagant. Pravda on January 8, for example, charged that \enquote{this organized racist campaign had been launched as if at a signal given by Adenauer} and it said that its purpose is to  \enquote{clear the road for a wide propaganda in favor of a nazi ideology, openly fascist.}})} the issue whether the most insidious form of hatred in Europe still persisted, possibly with subversive official influence, led to an inquiry by a UNGA subcommission, at the behest of Jewish organizations and American and Israeli political leaders.\footnote{Library Of Congress, and Sponsoring Body John W. Kluge Center, ``The Swastika Epidemic: Jewish Politics \& Human Rights in the 1960s,'' in Washington, D.C.: Library of Congress, -12-11, 2014. Video (transcript of lecture). \url{https://www.loc.gov/item/2021689625/}. (``In 1960  \ldots [as] news of these incidents [was published] across the world, a coalition began to form of international Jewish organizations, the Israeli Government and the U.S. officials. And, together, they launched an unprecedented effort to create a law that would ban anti-Semitism. They went to the United Nations, the Commission on Human Rights, seeking redress. And, this campaign, eventually, bore fruit.'')} 

In the process of deliberating on its recommendations, and at the prodding of some African member states,\footnote{Steven L. B. Jensen, \enquote{Lost Covenant: A Story of the Failed 1967 Convention on Elimination of Religious Intolerance,} January 12, 2017, Universal Rights Group. \url{https://www.universal-rights.org/lost-covenant-story-failed-1967-convention-elimination-religious-intolerance/}.}  the UNGA's Third Committee,\footnote{UNGA — Eighteenth Session, Resolutions Adopted On the Reports of the Third Committee. https://docs.un.org/en/A/res/1904(XVIII). See also, Eshed 2025, p. 249. On the UNGA's ``main committee'' structure, \textit{see generally}, Permanent Mission of Switzerland to the United Nations (2017). \textit{The GA Handbook: A practical guide to the United Nations General Assembly}. \url{https://unitar.org/sites/default/files/media/publication/doc/un_pga_new_handbook_0.pdf}.}  which oversees all social, humanitarian and cultural issues,\footnote{UN General Assembly - Third Committee. Social, Humanitarian \& Cultural, \url{https://www.un.org/en/ga/third/}.} split racial discrimination and religious intolerance into separate issues and proposed that resolutions be passed on each \parencite[257]{Thornberry-1991}. The UNGA duly accepted this without vote and directives were issued to prepare all such drafts in the interest of dealing with \enquote{manifestations of discrimination based on differences of race, colour and religion still in evidence throughout the world.}\footnote{UNGA Resolutions 1780 and 1781 (XVII). \url{https://docs.un.org/en/A/RES/1780 (XVII)}, \url{https://docs.un.org/en/A/RES/1781 (XVII}.} This led to the UNGA's November 1963 passage of a declaration on racial discrimination (Resolution 1904) followed by the ICERD in December 1965. But establishing consensus on the subject of religious intolerance proved far more difficult and the UNGA made no real progress until it finally passed a separate, relatively innocuous declaration  (because of the \enquote{conspicuous absence of incitement language}) in 1981.\footnote{Heiner Bielefeldt and Michael Wiener, \enquote{Declaration on the Elimination of All Forms of Intolerance and of Discrimination Based on Religion or Belief --- General Assembly resolution 36/55,} United Nations Audiovisual Library of International Law, 2021. 2. \url{https://legal.un.org/avl/pdf/ha/ga_36-55/ga_36-55_e.pdf}, 5.} As for the earlier draft, \enquote{the verbal confrontations during the debate, especially between Israel and the Arab states and India and Pakistan, reflected the extent to which ideological factors had intruded} in differing Muslim, Jewish and Hindu viewpoints.\footnote{John Claydon, ``The Treaty Protection of Religious Rights: U.N. Draft Convention on the Elimination of All Forms of Intolerance and of Discrimination Based on Religion or Belief,'' 12 Santa Clara Lawyer 403, 413 (1972). \url{https://digitalcommons.law.scu.edu/lawreview/vol12/iss2/8}.}   

As for the ICERD --- which was the first human rights convention the UNGA had adopted since 1948 --- its objectives were placed at the core of its \enquote{social, cultural and humanitarian} agenda going forward.\footnote{Steven L. B. Jensen, ``Lost Covenant: A Story of the Failed 1967 Convention on Elimination of Religious Intolerance,'' 2017, Universal Rights Group, \url{https://www.universal-rights.org/lost-covenant-story-failed-1967-convention-elimination-religious-intolerance/}. (``It is a common misinterpretation to view the post-1945 human rights story as driven by an emphasis on civil, political, economic, social and cultural rights. In reality, the international human rights project was redefined in 1962 with racial discrimination and religious intolerance at its core.'')} But before it was finalized, and after the United States and Brazil proposed adding a new article which identified anti-Semitism as a type of racism covered by the convention, the Soviet Union used that opportunity to advance its own ideological and geopolitical objectives.\footnote{The proposed revision by the United States and Brazil and the counter-revision by the Soviet Union appears on p. 5 of the Report of the Third Committee, UNGA (20th sess:), Agenda Item 53, Document A/6181 (December 18, 1965), \url{https://digitallibrary.un.org/record/712630?ln=en\&v=pdf}.} Ideologically, its concern was whether such an article could be used as the basis for challenging its own suppression of communal expressions of Jewish identity and Zionism in the Soviet Union and the satellite states it controlled after World War II because, following a brief hiatus that began soon after Stalin's death in 1953, such practices had resumed to full effect by the end of the 1950s. Geopolitically, it saw an opportunity to build a bridge to the Arab states that were no less furious than ever over Israel's intrusion into what had long been a solidly \enquote{pan-Arab} region.

This chain of events requires an understanding of shifting Soviet attitudes toward Israel, Jews and Zionism from the late 1940s to the late 1950s. While relatively relaxed up to the first half of 1948 \parencite[234]{Pinkus-1984}, when the USSR supported the post-mandate partition of Palestine --- based largely on Stalin's belief that Israel's socialist-leaning leaders would prove amenable to Soviet interests --- relations had soured significantly by 1949-50 due to several factors. These included Stalin's alarm over a surprisingly massive Rosh Hashanah celebration in Moscow that coincided with a visit by Israeli ambassador Golda Meir in the fall of 1948 (reflecting Russian Jewry's suspicious enthusiasm for the new Jewish state), as well as Israel siding with the United States' position at the UN regarding military intervention on the Korean peninsula.\footnote{Palquest, \url{https://www.palquest.org/en/highlight/38300/ussr-and-palestine-question-1950\%E2\%80\%931991}.} Thus, for example, at the behest of their Soviet handlers, the staff of the leading Jewish-American Communist newspaper unexpectedly declared in 1949 that Zionism was ``Jewish bourgeois nationalism'' and Israel, ``a bastion of Jewish reaction,'' served as ``a prop for Anglo-American imperialism.''\footnote{Harry Schwartz, \enquote{Communists Here to Fight Zionism; Freiheit Association Apologizes for Not Embracing Stand of Ehrenburg Earlier}, \textit{New York Times}, March 31, 1949. \url{https://www.nytimes.com/1949/03/31/archives/communists-here-to-fight-zionism-freiheit-association-apologizes.html}.} At the same time, Stalin grew wary of a nascent reawakening of Jewish organizations and activities in Eastern Europe which appeared \enquote{preoccupied with Jewish history} and inclined to slander \enquote{Soviet man and Soviet reality} \parencite[148--149]{Pinkus-1984}, making it suspect as a servant of \enquote{international Zionism and part of a conspiracy opposing socialism.}\footnote{Izabella Tabarovsky, ``Soviet Anti-Zionism and Contemporary Left Antisemitism'', \textit{Fathom Journal}, May 2017. \url{https://fathomjournal.org/soviet-anti-zionism-and-contemporary-left-antisemitism}.} Renewed vilification ended soon after Stalin's death, which appears attributable to ``liberalization processes'' associated with the Yugoslavian leader Josep Tito's influence on Khrushchev to accept the ``inevitability of a showdown with Stalin's cult.''\footnote{Svetozar Rajak, ``Yugoslav-soviet Relations, 1953-1957: Normalization, Comradeship, Confrontation' (Ph.D. thesis),' 2004, London School of Economics and Political Science, p. 343, \url{https://researchonline.lse.ac.uk/id/eprint/133506/}.} But it resumed a few years later, when Yiddish books and theater again proliferated as a result of such processes, leading to the resumption of \enquote{fierce campaigns against the Jewish religion and Zionism} \parencite[151]{Pinkus-1984}.

It is against this backdrop that Soviet actions with respect to final modifications of the ICERD's wording should be understood. While the additional article proposed by the U.S. and Brazil would confirm that anti-Semitism was within the purview of the ICERD, the Soviet counter-proposal --- offered contingent on that article's inclusion --- was provocatively designed to protect its own parochial priorities. Faced with a counter-proposal that the new article make clear that opposition to Zionism would not be affected, because it should be equated with ``Nazism, neo-Nazism and all other forms of the policy and ideology of colonialism, national and race hatred,'' the commission tasked with submitting the final version of the ICERD to the UNGA for a vote opted instead to drop both the American proposal and Soviet counter-proposal. But importantly, in preventing a reference to antisemitism from being explicitly inserted into the ICERD, the Soviets also took a stance that was largely consistent with common understandings of the era, that antisemitism was not a form of racism because the Jews could not be considered a `race.'  


The key to this puzzle is similar to the quandary often pondered today regarding the recent resurgence of anti-Zionism when juxtaposed against muddled understandings of the precise nature of antisemitism, as generally formulated in the mid-twentieth century and continuing for decades thereafter. Although it is widely recognized (at least in certain ways) as a coherent whole, the world's Jewish population was commonly perceived as a ``mixed group racially'' \parencite[268]{Graeber-1953}. This manifested widely-held assumptions that even while a ``scientific'' concept of race must be discredited as a factor in fundamental behavioral differences between population groups, because there really were none, the concept still had some validity and deserved credence. It could serve as an overarching classification system based on physical (if not intellectual or emotional) differences within humanity's shared gene pool, broadly associated with different traits, whether resulting from a mix of genetic inheritances and environmental factors or otherwise. Such differences, whether superficial or in some ways possibly substantive, were considered easily detectable between diverse populations originating from regionalized geographies. 

The amount of scholarly literature on this complex topic --- involving human biology, anthropology and sociology --- has been vast for well over a century, and it would be foolish to declare any single point-of-view about it as the correct one. But the broader problem lies in the fact that by renouncing the credibility of \enquote{scientific racism,} because it was impossible to verify in any real way through newly discovered genetic methodologies, scholars in fields of study that are less amenable to empirical data --- philosophy and the social sciences, and particularly the relatively new-fangled discipline of sociology --- were in search of workable concepts that explained the persistence of conflict among separately identifiable population groups. Thus, while it became obvious that \enquote{everyday racial classifications do not track meaningful biological categories,}\footnote{Sally Haslanger, ``You mixed? Racial identity without racial biology,'' In: S. Haslanger S \& C. Witt (eds), \textit{Adoption matters: philosophical and feminist essays}, 266. (Cornell University Press: Ithaca, 2005)} the concept of race still demanded study because of its persistent and often deeply disruptive sociological impacts, even though now considered devoid of, and therefore independent of, strictly biological determinants. 

Without widely accepted and certainly in no way ``fixed'' definitions of race, racism or antisemitism, hatred toward Jews could therefore be placed within a category adjacent to conflict attributable to race, consisting of extreme ethno-religious intolerance, but not squarely within it, with race still largely understood as relating to the ``immutable'' characteristics of distinct, region-specific populations, reflecting hereditary continuities and ``genetically informed'' --- if not decisive --- factors \parencite[160]{Sesardic-2010}. Antisemitism could thus be categorized as hatred toward the presence of a transnational minority demarcated by unique ethnic identifiers, such as ``wearing distinctive garb'' in their everyday lives or ``characteristically Jewish names'' \parencite[268]{Graeber-1953}. But this hatred was  not ``racial'' because it could not be considered strongly tied to physical characteristics such as skin color, facial structure, hair texture, which could be associated with other groups --- whether accurately or imaginatively --- such as was the case with \enquote{white Caucasians} whose origins were found in an imaginary yet still specific geography.\footnote{A recent book addressing the construction of a ``Caucasian'' race and its associated myths by Sarah Lewis, \textit{The Unseen Truth: When Race Changed Sight in America}, Harvard University. Press: 2024) explores this phenomenon from the perspective of the homogenizing \enquote{racialization} of a multiplicity of otherwise distinct in many ways (culturally, linguistically, and physically) Eurasian  peoples.} The point is relevant here because the rigid use of physical factors to define racial differences only became fully discredited over the course of the second half of the last century and, while the qualities that define what exactly is a \enquote{race} have either been dispensed with or grown less pseudo-scientific as a result, they are still relatively murky in terms of shared understandings, or even widespread acceptance, today.\footnote{E.g., \enquote{Is Antisemitism a Form of Racism?} \textit{Anne Frank House}, (undated).  \url{https://www.annefrank.org/en/topics/antisemitism/antisemitism-form-racism/} (``Antisemitism means hatred of Jews. The word first appeared in the 19th century, when classification of people into different races was considered normal \ldots. But do Jews belong to a separate `race'? \ldots. Conclusion: Jews are not a race, and categorising people according to race is wrong and dangerous. Even so, some people still believe in the concept. If it is the basis for their hatred of Jews, it is undoubtedly racist.'')} Thus, a unifying hatred toward the Jews, even now, might be seen as an integral part of the ``gross zoological racism and anti-semitism'' of the Nazis, but the two hatreds --- racism and antisemitism --- could still be recognized and treated as complementary but not one and the same.\footnote{Aurel Kolnai, \textit{The War Against the West}, Viking Press 1938, 38. Also, regarding the ideas of \enquote{Nordicism} --- that race was the basis of civilization [with the assumption] that it must therefore ``be the basis of political order,'' \textit{see}, John P. Jackson and Nadina M. Weidman, \enquote{The Origins of Scientific Racism,} The Journal of Blacks in Higher Education (2005-05), 66--79, 50. \url{https://www.jstor.org/stable/25073379}.} 

A precise chronology of such academic reconfigurations, or even an assessment of their relevance today, is all but impossible to arrive at --- even though attempts to do so are still made with respect to the connection between antisemitism and racism\footnote{Glynis Cousin \& Robert Fine (2012) ``A Common Cause: Reconnecting the study of racism and antisemitism,'' European Societies, 14:2, 166-185. \url{https://doi.org/10.1080/14616696.2012.676447}} --- in part because the terms, concepts and meanings in play are all inherently transitory and opaque. Nevertheless, the notion that \enquote{racial identity} can be formed even if not \enquote{based on attributes typically linked to physical characteristics like skin color,} thereby entitling members of a population with diverse physical characteristics to consider themselves belonging to a \enquote{racialized} group by virtue of \enquote{ascription (how society labels an individual), identification (how one self-identifies), and treatment (how one is treated by others),}\footnote{Asil M. Martinez, \enquote{Is Racial Identity Ever Straightforward?}, Blog of the APA, June 10, 2024. \url{https://blog.apaonline.org/2024/06/10/is-racial-identity-ever-straightforward/}.} has become a bulwark of modern thinking. This has allowed nearly endless and often perambulating explorations on what groups have been (or are in the process of being) racialized, what indicates their racialization, and what are the implications of their having been racialized for themselves and others.

Somewhat amazingly, this flexible, case-sensitive and highly erudite area of inquiry now lies at the heart of international law with respect to what groups could be considered subject to a regime constituting apartheid. The non-scientific acceptance  of this outcome is detectable not only in the copious academic literature on the concept of \enquote{racialization} in the last several decades, but also the deployment of the term \enquote*{racial group,} superseding that of \enquote*{race,} that now permeates the wording of international conventions. It began with the ICERD, where the term appeared three times in an undefined, relatively off-handed manner\footnote{Art. 1, sec. 4 addresses \enquote{the maintenance of separate rights for different racial groups} not being continued after special measures are taken to secure the \enquote{adequate advancement of certain racial or ethnic groups.}  protection of certain racial groups or individuals belonging to them; and Art. 2, sec. 2 uses the term twice in relation to the taking of ``concrete measures'' to ensure \enquote{adequate development and protection of certain racial groups or individuals belonging to them}.} and became more apparent with the UNGA's 1974 \enquote{Apartheid Convention,} where it was used 13 times, including in the definition of ``the crime of apartheid.'' It was still undefined, but it thus became integral to the crime's provability.\footnote{International Convention on the Suppression and Punishment of the Crime of Apartheid, UNGA Res. 3068 (XXVIII), 28 U.N. GAOR Supp. (No. 30) at 75, U.N. Doc. A/9030 (1974), 1015 U.N.T.S. 243, entered into force July 18, 1976. \url{https://www.un.org/en/genocideprevention/documents/atrocity-crimes/Doc.10_International\%20Convention\%20on\%20the\%20Suppression\%20and\%20Punishment\%20of\%20the\%20Crime\%20of\%20Apartheid.pdf}.} According to defenders of an expansive approach toward utilizing racialization as an analytical process, it is of \enquote{great value politically, as it offers a way for groups that have been understood and treated as inferior \enquote{races} to assert and defend themselves collectively, while rejecting the biologization and inferiorization associated with \enquote{race.}}\footnote{Adam Hochman, ``Racialization: a defense of the concept,'' Ethnic and Racial Studies, 42(8), 1245--1262, 1245-46. \url{https://doi.org/10.1080/01419870.2018.1527937}.} The process of identifying racial groups since the 1970s has also become something of a thriving cottage industry in social sciences academia.\footnote{Steve Garner, \textit{Racisms: An Introduction.} SAGE Publications Ltd, 2010), 19. (``The concept of racialisation is based on the idea that the object of study should not be \enquote*{race} itself, but the process by which it becomes meaningful in a particular context. In fact, racialisation has now become one of the key ways that academics make sense of the \enquote*{meanings of race}.''} Thus, while the physiological foundations of race-based differences have persisted --- often denoted by skin color (leading to the now widely-used concept of ``colorism'' as a sub-type of racism)\footnote{Martin Lund (September 10, 2024), ``A Prehistory of Scientific Racism.'' The MIT Press Reader. \url{https://thereader.mitpress.mit.edu/a-prehistory-of-scientific-racism/}. (``In the 1795 third edition of his pamphlet \enquote*{On the Natural Variety of Mankind,} German anthropologist Johann Friedrich Blumenbach \ldots skin color played a large role in what was considered the science of race. He included skin color in his taxonomy and ranked white skin the highest, as belonging to the \enquote*{oldest variety of man,} although skin color was ascribed to climate and experience rather than innate qualities, and \enquote*{Caucasian} whiteness was extended as far as the Ural Mountains and Ganges.'')} --- racism (or more precisely, ``racial bias'' toward a ``racial group'') is now susceptible to understanding and interpretation with as much malleability as deemed necessary to make it considered fair and equitable, and therefore susceptible to the imperatives of international law.\footnote{An example of the trend toward empirical evaluation of racial bias as a social phenomenon can be found in R. Neel and J.R. Shapiro, ``Is racial bias malleable? Whites' lay theories of racial bias predict divergent strategies for interracial interactions,'' \textit{Journal of Personality and Social Psychology}, May 7, 2012, 103(1), 101--120. \url{https: https://pubmed.ncbi.nlm.nih.gov/22564011/}.} 

This broad and in some ways obscure topic may seem largely separable from the history of the UNGA's affairs being examined here because, throughout the duration of South Africa's long apartheid era, there would have been no obvious need for its use, much less close examination. That is, from the 1960s through at least the early 1990s, widely-accepted identifiers of different populations based on established concepts of racial difference did not require sociological analysis or nuanced theory. Except for a relative handful of edge cases, the outward and obviously manifested basis for all-but-unquestioned separate and unequal treatment between \enquote{Whites} and \enquote{`Coloureds} would have been plainly obvious. Yet when the UNGA moved on from cursory references to racial groups in the ICERD to adopting the Apartheid convention 10 years later, the purpose of \enquote{racial group} as the numerator to the denominator of the rest of the population in question became both more versatile, and more entrenched as particularly meaningful, for purposes of scope, interpretation and effect. 

First proposed by the USSR, Nigeria and Guinea\footnote{Newell M. Stultz, \enquote{Evolution of the United Nations Anti-Apartheid Regime}, 13 Human Rights Quarterly  1--23 (1991), 13} and submitted in Russian,\footnote{Guinea and Union of Soviet Socialist Republics: draft of a Convention on the suppression and punishment of the crime of apartheid (28 Oct. 1971) UNGA AC.3/L.1871.} the Apartheid convention was seen as \enquote{the ultimate step in the condemnation of apartheid as it not only declared that apartheid was unlawful because it violated the Charter of the United Nations, but in addition it declared apartheid to be criminal.}\footnote{John Dugard, \enquote{Introductory Note, Convention on the Suppression and Punishment of the Crime of Apartheid,UN Audio Visual Library of International Law, August 2008. \url{https://legal.un.org/avl/ha/cspca/cspca.html}.}} However, given the plethora of actions that had been continually taken in the UNGA and by the security council regarding apartheid up to that point, member-state submissions  \enquote{revealed a great variety of opinions on the need for the draft convention as such}\footnote{S. Azadon Tiewul, \enquote{Apartheid: Steps Toward International Legal Control, 119} 6 ZAM. Law Journal 101, 119 (174)} and suspicions were raised that the convention was more a symbolic gesture of ongoing Soviet solidarity with newly independent, anti-imperialist African states than a practical instrument. It included a monitoring body but appointed members on it needed to be \enquote{neither independent nor expert.}\footnote{Victor Kattan \& David Johnson, \enquote{The Crime of Apartheid beyond Southern Africa: A Call to Revive the Apartheid Convention's \enquote{Group of Three},} EJIL: Talk!, Blog of the European Journal of International Law, September 21, 2023, \url{https://www.ejiltalk.org/the-crime-of-apartheid-beyond-southern-africa-a-call-to-revive-the-apartheid-conventions-group-of-three/}. See also, Dugard 2008 (``Although consideration was given in 1980 to the establishment of a special international criminal court to try persons for the crime of apartheid (E/CN.4/1426 (1981)), no such court was established.'')} It was \enquote{intended to complement the ICERD}\footnote{Max Du Plessis, \enquote{International Criminal Law: The Crime of Apartheid Revisited,} 24 S. African Journal of Criminal Justice. 417--428, 4220 (2011).} but the crime it defined and declared came without any ways to prosecute acts that constituted it, or enforce penalties for doing so.  

Perhaps most importantly, while the Apartheid convention identified eight acts demonstrative of the ``purpose of establishing and maintaining domination by one racial group of persons over any other racial group of persons and of systematically oppressing them,'' it referenced them as ``''\textit{including}'' those \textit{``as practised in southern Africa.}''\footnote{Apartheid Convention, Art. I, sec. 2: \enquote{[T]he term `the crime of apartheid', which shall include similar policies and practices of racial segregation and discrimination as practised in southern Africa, shall apply to the following inhuman acts committed for the purpose of establishing and maintaining domination by one racial group of persons over any other racial group of persons and systematically oppressing them \ldots.}} Expert legal commentators subsequently noted that this broadened the convention's potential application to political and social conditions far beyond South Africa, regardless of the extent to which they were like or unlike it \parencite[211]{Reynolds-2012}. Through this apparent reference to South Africa's case as an example but not a set of parameters,   apartheid shifted from a subcontinental crisis resulting from an array of specific circumstances --- the multi-century development of Afrikaner history and culture; a national economy heavily dependent on mineral extraction; the milieu of an originally isolated outpost in a lightly populated area meant to facilitate the commercial expansion of Dutch mercantilism \parencite [xv--xviii]{Brookes-2023}; and a post-World War II climate of aggressive White supremacist nationalism to protect at all costs the material gains of what had long been an environmentally precarious existence --- was thereupon imbued with potentially global application.

In any case, and over the long run, the Apartheid convention's potential was largely theoretical, and it is now \enquote{widely believed to be a dead letter} \parencite[86]{Lingaas-2015}, given that its monitoring function ended in 1995, once South Africa became a post-apartheid state.\footnote{Kattan \& Johnson 2023.} But it set the stage for a newly internationalized crime against humanity whose many champions --- now in a worldwide search for possible perpetrators --- live on. And, like the ICERD before it, the convention's adoption in itself reflected a victory, whether pragmatic or ideological. For the post-colonial states in Asia and Africa, it showed them that their Communist partners were willing to protect their interests, even if their primary reasons for doing so were to combat western influence \parencite[335--338]{Golan-1988} and maintain ideologically allied regimes in Angola and Mozambique.\footnote{CIA Directorate of Intelligence, \enquote{Soviet Intentions and Activities in Southern Africa,} November 25, 1986 \url{https://www.cia.gov/readingroom/docs/CIA-RDP86T01017R000505430001-3.pdf}, 2.} The strategy clearly worked because the vast majority of UNGA member-states voted in favor of the convention\footnote{Multilateral Treaties in respect of which The SECRETARY-GENERAL Peforms Depositary Functions --- List of Signatures, Ratifications, Accessions, etc. as at 31 December 1974, 102. \url{https://treaties.un.org/doc/source/publications/MTDSG/1974-english.pdf}.} and only a relative handful either opposed or abstained.\footnote{Dugard 2008.} 

Of equal if not greater significance to the value of the Apartheid convention in ``instrumentalizing [the USSR's] relations with the Global South,'' it helped characterize apartheid as reflective of the natural tendencies of western imperialism and colonialism, with the result being that it solidified the loyalties of ``non-aligned movement'' countries during the Cold War years even while they kept their diplomatic distance from the Soviets' more affluent Western adversaries, no matter how helpful they might be in their offers of aid for much needed economic development.\footnote{Drachewych, O. (4 Mar. 2025) The Comintern and the National and Colonial Question: Reconsidering the Roots of Soviet Anti-Imperialism and Anti-Racism. The Jordan Center for the Advanced Study of Russia at New York University.  \url{https://jordanrussiacenter.org/blog/the-comintern-and-the-national-and-colonial-question-reconsidering-the-roots-of-soviet-anti-imperialism-and-anti-racism}. (``Through the Comintern, the communist movement developed many foundational principles of what would become Soviet anti-imperialism and anti-racism \ldots. It is possible that Soviet support for anti-apartheid or the introduction of leftist ideas into civil rights movements in the United States would not have taken place without the forum the Comintern offered.'') The author is a Canadian academic who specializes in ``the history of the Soviet Union, international communism, and global history and international relations in the twentieth century.'' \url{https://www.kings.uwo.ca/academics/faculty-info/member-profile/?doaction=getProfile\&id=odrachew}.} This strategy can be traced back not just to ten years earlier, with the 1965 adoption of the ICERD, but 20 years, to the 1955 Bandung Conference, convened ``to explore ways that newly independent states could navigate the Cold War bipolarity between a U.S.-led `capitalist' bloc and the Soviet-led `socialist' bloc.''\footnote{B.N. Mamlyuk, ``Decolonization as a Cold War Imperative: Early Soviet International Law as a Precursor to Bandung,'' SSRN Electronic Journal, 2016. \url{https://papers.ssrn.com/sol3/papers.cfm?abstract_id=2601524}. This article is a ``preprint'' version of a chapter in a 2016 book cited in the next footnote, which (for unknown reasons) does not appear to make precisely the same point.} In fact, Bandung may be considered the real watershed start of the era because, in its success in bringing together \enquote{leaders of countries whose combined population made up approximately two-thirds of the world's people,}\footnote{Luis Eslava, Michael Fakhri, and Vasuki Nesiah, \enquote{The Sprit of Bandung,} 3--32, 4, in \textit{Bandung, Global History, and International Law: Critical Pasts and Pending Futures,} Luis Eslava, Michael Fakhri, Vasuki Nesiah (eds.), Cambridge University Press 2017.} it garnered worldwide attention and caused \enquote{Soviet theory} of its overall engagement strategy for the Cold War to change dramatically, from Stalin's \enquote{two-camp world} to Khrushchev's third \enquote{zone of peace,} the recognition of which led the Soviets to embark on \enquote{a period of winning friends and influence amongst the newly independent former colonies.}\footnote{Mayluk, 81.}

The cumulative effect of such events during this quarter-century era (from roughly 1950 to 1975) is now abundantly clear a half century after it ended: the struggle against racism, especially in its most brutal, pernicious and state-sanctioned forms, must be directed against Western habits --- indeed, requirements --- of developing \enquote{racial group} identities that are brought into existence for the sole purpose of subjugating them to achieve imperialist objectives. The UNGA that emerged after South Africa's apartheid system collapsed, heavily influenced throughout this era by the durable coalition of member states that purveyed this messaging, might be described as the organizational equivalent of a St. George dedicated to achieving a type of victory which it can no longer redefine and from which it can no longer retire.\footnote{Cf., Kenny Minogue, \textit{The Liberal Mind,} Vintage Books, 1968, 1.} It managed to slay the primary dragon of the last century --- South Africa's apartheid regime --- but in the new century, what other dragon-like regimes, inflicting immense harms on otherwise helpless \enquote{racial groups,} must be recognized, shunned and slain now? As it turns out, there was at least one that was readily available.

\clearpage
\parindent 0pt
\parskip 6pt

\end{spacing}

\setenotez{
  list-name = {Notes to Chapter \thechapter}
}
% Place this macro at the exact end of every chapter:
{\footnotesize
\printendnotes[paragraph]}

\thispagestyle{empty} % makes a blank page

\end{document}